% This command tells LaTeX what kind of document we are creating (article).
\documentclass{article}

\usepackage{amsmath}
\usepackage{array}
\usepackage{booktabs}
\usepackage[margin=1in]{geometry}
\usepackage{graphicx}
\usepackage{lipsum}
\usepackage[parfill]{parskip}
\usepackage{xcolor}

% Highlight Keywords using the \kw{} command
\newcommand{\kw}[1]{\textcolor{blue}{\textbf{#1}}}
% Italicise LaTeX commands
\newcommand{\cmd}[1]{\textit{#1}}

% Everything before the \begin{document} command is called the preamble.
\begin{document} % The document body starts here

\begin{titlepage}
    \centering

    \huge
    \textbf{My Example Document}

    \vspace{1cm}
    \normalsize
    \textit{An example of a LaTeX document}

    \vfill
    January 1, 2000
\end{titlepage}

Hello World!

This is my first \kw{LaTeX} document.

% The section command automatically numbers and formats the section heading.
\section{Sections}

I can add content to my first \kw{section}!

% The subsection command does the same thing, but for sections within sections.
\subsection{Subsection}

I can put a \kw{subsection} inside my first section.

\section{Lists}

There are two types of \kw{lists}: \kw{ordered} and \kw{unordered}.

\subsection{Ordered}

Ordered lists do not have numbers associated with each item.

\begin{enumerate}
  \item Item 1
  \item Item 2
  \item Item 3
\end{enumerate}

\subsection{Unordered}

Unordered lists are just a series of items preceded by a marker.

\begin{itemize}
  \item Item 1
  \item Item 2
  \item Item 3
\end{itemize}

\section{Graphics}

We can include \kw{images} in our document using the \cmd{graphicx} package, which lets us use the
\cmd{includegraphics} command.

\includegraphics{example-image.png}

\subsection{Small Image}

We can pass parameters to the \cmd{includegraphics} command to adjust the appearance of the image.

\includegraphics[height=2cm]{example-image.PNG}

Other possible options include:

\begin{itemize}
  \item width: the width of the image
  \item scale: the scaling factor of the image
  \item angle: the angle of rotation of the image
  \item clip: whether to clip the image to its bounding box
  \item trim: trim the image by a specified amount
  \item draft: display a box instead of the image
\end{itemize}

\subsection{Centered Image}

By placing the \cmd{includegraphics} command inside a center environment, we can center the
image on the page.

\begin{center}
  \includegraphics[height=2cm]{example-image.PNG}
\end{center}

\subsection{Floating Image}

\kw{Floating images} can move around the page as text is added or removed. We can use the
\cmd{figure}environment to create a floating image.

\begin{figure}[ht]
  \centering
  \includegraphics[height=2cm]{example-image.PNG}
\end{figure}

Control the position of a floating image by passing paratmeters to the \cmd{figure} environment:

\begin{itemize}
  \item h: Place the float "here" (where it appears in the code)
  \item t: Place the float at the "top" of the page
  \item b: Place the float at the "bottom" of the page
  \item p: Place the float on a "page" by itself
\end{itemize}

\subsection{Caption}

We can add a \kw{caption} to our floating image by using the \cmd{caption} command inside of the
\cmd{figure} environment.

\begin{figure}[h]
  \centering
  \includegraphics[height=2cm]{example-image.PNG}
  \caption{This is a caption for our image.}
\end{figure}

\section{Tables}

\kw{Tables} are defined using the \cmd{tabular} environment.

\begin{tabular}{*{3}{l}}
  Fruit  & Quantity & Price \\
  Apple  & 5        & 1.50  \\
  Banana & 6        & 2.00  \\
  Orange & 4        & 1.20  \\
\end{tabular}

\subsection{Tables with Horizontal Lines}

We can use the \cmd{\textbackslash toprule}, \cmd{\textbackslash midrule}, and
\cmd{\textbackslash bottomrule} commands from the \cmd{booktabs} package to create horizontal
lines in our table.

\begin{tabular}{*{3}{l}}
  \toprule
  Fruit  & Quantity & Price \\
  \midrule
  Apple  & 5        & 1.50  \\
  Banana & 6        & 2.00  \\
  Orange & 4        & 1.20  \\
  \bottomrule
\end{tabular}

\subsection{Partial Horizontal Lines}

The \cmd{\textbackslash cmidrule} command can be used to create partial horizontal lines in a
table. The command accepts the arguments {number-number}, where the first number is the column to
start the line and last number is the column to end the line.

\begin{tabular}{*{3}{l}}
  \toprule
  Fruit  & Quantity &  Price \\
  \midrule
  Apple  & 5        &  1.50  \\
  Banana & 6        &  2.00  \\
  Orange & 4        &  1.20  \\
  \cmidrule{3-3}
  Total & 15        & 28.20  \\
  \bottomrule
\end{tabular}

\subsection{Merging Cells}

Merge cells horizontally using the \cmd{\textbackslash multicolumn} command. This command takes
three arguments:

\begin{itemize}
  \item The number of cells which should be merged
  \item The alignment of the merged cell (l, c, or r)
  \item The contents of the merged cell
\end{itemize}

\begin{tabular}{*{3}{l}}
  \toprule
  \multicolumn{3}{c}{Overall Inventory} \\
  Fruit  & Quantity &  Price  \\
  \midrule
  Apple  & 5        &  1.50   \\
  Banana & 6        &  2.00   \\
  Orange & 4        &  1.20   \\
  \midrule
  \multicolumn{2}{c}{Summary} \\
  Total & 15        & 28.20   \\
  \bottomrule
\end{tabular}

\section{Cross References}
\label{sec:cross-references}

\subsection{Material for the Introduction}

In this section, we introduce two new concepts:

\begin{table}[ht] % need to add a table environment
  \begin{tabular}{cp{9cm}}
     Command & Description \\
     \toprule
     \kw{label} & Marks a spot in the document \\
     \kw{ref} & Refers to a marked spot in the document \\
     \bottomrule
  \end{tabular}
  \caption{This is needed for the label to work.}
  \label{tab:cross-reference-commands} % the lable refers to the caption
\end{table}

We can \kw{refer} to the section with the \cmd{ref} command, like this: \ref{sec:cross-references}.
We can likewise refer to the table like this: \ref{tab:cross-reference-commands}.

\section{Mathematics}

There are two kinds of math mode in LaTeX: inline and display. Inline math mode is marked with
a pair of dollar signs, whereas display math mode is marked with a pair of square brackets.

\subsection{Inline Math Mode}

The Pythagorean theorem is $a^2 + b^2 = c^2$.

\subsection{Display Math Mode}

The Fourier Transform is defined as:

\[
\hat{f}(\xi) = \int_{-\infty}^\infty f(x) e^{-2\pi i \xi x} \, dx
\]

Where:
\begin{itemize}
  \item \( f(x) \) is the function we are transforming,
  \item \( \hat{f}(\xi) \) is the Fourier Transform of \( f(x) \),
  \item \( \xi \) is the frequency variable,
  \item \( i \) is the imaginary unit.
\end{itemize}

\subsection{Math in Environments}

The quadratic formula is:

\begin{equation}
x = \frac{-b \pm \sqrt{b^2 - 4ac}}{2a}
\label{eq:quadratic}
\end{equation}

This will allow us to refer to the equation later in the document with \cmd{ref} like this:
Refer to Equation \ref{eq:quadratic}.

\subsection{The `amsmath` Package}

Solve the following recurrence for $ n,k\geq 0 $:
\[
Q_{n,0} = 1   \quad Q_{0,k} = [k=0];
\]

\[
Q_{n,k} = Q_{n-1,k}+Q_{n-1,k-1}+\binom{n}{k}, \quad\text{for $n$, $k>0$.}
\]

\section{Fonts and Spacing}

% Generate some "Lorem Ipsum" text
% The parameters mean "include paragraphs 1 thru 2" from the "Lorem Ipsum" text
\lipsum[1-2]

\end{document}
